Deux blocs posés sur une table horizontale sans frottement sont reliés par une corde de masse négligeable. Le bloc 1 a une masse de \qty{\v{m1}}{kg} et le bloc 2, de \qty{\v{m2}}{kg}. On tire horizontalement sur le bloc 2 avec une force de \qty{\v{F}}{N}, de sorte que le bloc 1 suit derrière.

\begin{enumerate}[label=\alph*)]
    \item Faites le diagramme des forces de chaque bloc et déterminez l'accélération du système.
    \item Déterminez la tension dans la corde qui relie les blocs.
    \item Partant du repos, quelle distance les blocs parcourent-ils en \qty{\v{t}}{s} ?
\end{enumerate}
